\input{config/documento}

\title{Título do Trabalho}

\def\autorA{Nome do Aluno 1\\
\texttt{meu.email@academico.ufpb.br}}

\def\autorB{Nome do Aluno 2\\
\texttt{meu.email@academico.ufpb.br}}

\def\disciplina{Técnicas de Pesquisa e Análise de Dados I}
\def\professor{Aléssio Tony C. Almeida}


\begin{document}
\input{config/pretextual}

\chapter{Introdução}

%Texto aleatório
\lipsum[1-1]
	
	%Exemplos de citação:
	
	\begin{itemize}
		\item De acordo com \cite{Almeida2019} e \cite{Silva2019}, o Brasil...
		\item O Programa Farmácia Popular do Brasil visa aumentar o acesso a medicamentos para doentes crônicos no país, usando especialmente a rede conveniada ao programa \citep{Almeida2019}.
	\end{itemize}
	
%Texto aleatório
\lipsum[1-5]

\section{Objetivos}
O objetivo geral deste trabalho é analisar......

\subsection{Objetivos Específicos}

\begin{itemize}
    \item Analisar....
    \item Identificar...
    \item Avaliar...
\end{itemize}

\chapter{Metodologia}

Descrever nesta parte a metodologia do seu trabalho.

\section{Método}
Descrever o(s) método(s) de análise de dados utilizado(s).

\section{Dados}
Descrever as variáveis utilizadas, período de análise, abrangência territorial do estudo, tamanho da amostra e as fontes de dados, bem como apresentar uma tabela com uma sumarização estatística das variáveis usadas (média, desvio-padrão, mínimo, máximo).

\chapter{Resultados}
Apresentar os achados do trabalho a fim de responder o problema fundamental de pesquisa.

\chapter{Conclusões}
Informar neste capítulo as principais considerações da pesquisa: resultados-chave, limitações, perspectivas futuras de trabalhos relacionados ao tema...

\bibliography{referencias}
\end{document}